\documentclass[11pt]{article}

\usepackage{amsmath,amssymb}
\usepackage{xurl}
\usepackage[hidelinks]{hyperref}
\setlength{\emergencystretch}{2em}

\input{command}

% Local notation matching the source paper (its style file is not shared here).
\newcommand{\val}{\operatorname{val}}
\newcommand{\game}{\mathfrak{G}}
\newcommand{\strategy}{\mathcal{S}}

\title{Attainment in the Symmetrization Lemma}
\author{}
\date{2026-08-30}

\begin{document}
\maketitle

\section{At a glance}

\begin{itemize}
  \item \textbf{Difficulty: easy.}  The issue is a quantifier exchange between
  a supremum and an attained maximum in a preliminary lemma on nonlocal games.
  No analytic estimate is missing; the substance is deciding which statement
  the symmetrization proof supports and which statement the later arguments
  consume.
  \item \textbf{Estimated weight:} about \textbf{1--2 lemma interfaces} and on
  the order of \textbf{100--300 lines} of formal proof for the value-preserving
  symmetrization construction.  The attainment step itself should not be
  formalized as printed (see the Conclusion).
  \item \textbf{Mathlib/project split:} approximately \textbf{20\% Mathlib}
  (finite-dimensional inner product spaces, tensor products, block
  decompositions of $\C^2 \ot \C^d$) and \textbf{80\% project-local} (games,
  strategies, values, and the Naimark input).
  \item \textbf{Key mathematical inputs:} the completed projective form of
  Naimark dilation, \cite[Theorem~5.1]{Ji2020LowIndividualDegree}, and the
  definitions of symmetric games and tensor product values in the source
  \cite[Section~5]{Ji2020MIPStar}.  Tracked in \localissue{0002}.
\end{itemize}

\section{Key theorem forms}

Three statements must be kept apart.  Throughout, $\game$ denotes a symmetric
two-player one-round game and $\eps \ge 0$ a real parameter; the notation is
introduced in the next section.

\begin{enumerate}
  \item \textbf{Source form (attainment form).}  If
  $\val^*(\game) = 1-\eps$, then there exists a symmetric and projective
  strategy $\strategy$ with $\val^*(\game, \strategy) \ge 1-\eps$.
  \item \textbf{Approximate form (what the source proof establishes).}  If
  $\val^*(\game) = 1-\eps$, then for every $\eps' > \eps$ there exists a
  symmetric and projective strategy $\strategy$ with
  $\val^*(\game, \strategy) \ge 1-\eps'$.
  \item \textbf{Given-strategy form (value-preserving symmetrization).}  If
  $\strategy'$ is any strategy for $\game$ with
  $\val^*(\game, \strategy') \ge 1-\eps$, then there exists a symmetric and
  projective strategy $\strategy$ with
  $\val^*(\game, \strategy) = \val^*(\game, \strategy') \ge 1-\eps$.
\end{enumerate}

The source proof consists exactly of a proof of the third form, applied to a
strategy supplied by the definition of the supremum; this yields the second
form.  The first form asserts strictly more than the second, and the passage
from the second to the first is the gap discussed in this note.  Every use of
the lemma, in the source and in this project, consumes only the second or the
third form.

\section{The source assertion}

A two-player one-round game
$\game = (\mathcal X, \mathcal A, \mu, D)$ in compact symmetric notation
consists of a finite question alphabet $\mathcal X$, a finite answer alphabet
$\mathcal A$, a symmetric question distribution $\mu$ on
$\mathcal X \times \mathcal X$, and a decision predicate $D$ treating the two
players symmetrically.  A tensor product strategy
$\strategy = (\ket\psi, A, B)$ consists of a unit vector $\ket\psi$ in a
tensor product $\mathcal H_A \ot \mathcal H_B$ of \emph{finite-dimensional}
Hilbert spaces together with POVMs $\{A^x_a\}$ and $\{B^x_a\}$ for the two
players; its value $\val^*(\game, \strategy)$ is the acceptance probability,
and the value of the game is the supremum
\begin{equation*}
  \val^*(\game) = \sup_{\strategy} \val^*(\game, \strategy)
\end{equation*}
over all tensor product strategies.  A strategy is projective if all its
measurements are projective, and symmetric if the state is
permutation-invariant on $\mathcal H \ot \mathcal H$ and both players use the
same measurements.

The symmetrization lemma of \cite[Section~5]{Ji2020MIPStar} asserts:%
\footnote{The statement is at
\path{references/qpbt-paper/06_nonlocal_games_and_mipstar.tex}, lines~94--99,
and its proof at lines~101--132.  The blueprint statement restored to this form
is \path{lem:symmetric-strat} in
\path{blueprint/src/chapter/ch12_qpbt_games.tex}.}
if $\game$ is a symmetric game with $\val^*(\game) = 1-\eps$ for some
$\eps \ge 0$, then
\begin{equation*}
  \exists\, \strategy = (\ket{\psi}, M) \text{ symmetric and projective}
  \quad\text{with}\quad
  \val^*(\game, \strategy) \ge 1-\eps.
  \tag{\path{lem:symmetric-strat}}
\end{equation*}

Its proof runs as follows.  For every $\eps' > \eps$ the definition of the
supremum supplies a strategy $\strategy' = (\ket{\psi'}, A, B)$ with
$\val^*(\game, \strategy') \ge 1-\eps'$.  By Naimark dilation in the completed
projective form of \cite[Theorem~5.1]{Ji2020LowIndividualDegree} one may take
$\ket{\psi'} \in \C^d_{A'} \ot \C^d_{B'}$ for some integer $d$ with all
measurements projective.  One then forms the permutation-invariant state
\begin{equation*}
  \ket{\psi} = \tfrac{1}{\sqrt 2}
  \bigl( \ket{0}_A \ket{1}_B \ket{\psi'}_{A'B'}
       + \ket{1}_A \ket{0}_B \ket{\psi'_\tau}_{A'B'} \bigr)
\end{equation*}
on $(\C^2 \ot \C^d) \ot (\C^2 \ot \C^d)$, where $\ket{\psi'_\tau}$ permutes
the two players' registers, and the block-diagonal projective measurements
$M^x_a = \ket{0}\!\bra{0} \ot A^x_a + \ket{1}\!\bra{1} \ot B^x_a$.  Symmetry
of $D$ gives $\val^*(\game, (\ket\psi, M)) = \val^*(\game, \strategy')$.  The
construction is therefore \emph{value-preserving}: applied to any strategy it
produces a symmetric projective strategy of exactly the same value.  What the
proof establishes is thus the given-strategy form, hence the approximate form;
the printed conclusion is the attainment form.

\section{Why attainment does not follow}

The supremum defining $\val^*(\game)$ ranges over strategies on
finite-dimensional Hilbert spaces of unbounded dimension.  The strategies
$\strategy'_{\eps'}$ produced for $\eps' \downarrow \eps$ may require
dimensions $d(\eps')$ growing without bound, and there is no compactness
across dimensions from which a limiting finite-dimensional strategy of value
at least $1-\eps$ could be extracted.  A symmetric projective strategy of
value at least $1-\eps = \val^*(\game)$ would attain the supremum, so the
printed conclusion is equivalent, under its hypothesis, to the assertion that
$\val^*(\game)$ is attained by some strategy.  No general principle supplies
this: it is known, and recalled in the introduction of the source itself, that
the set of finite-dimensional quantum correlations is not closed and differs
from the set of infinite-dimensional spatial correlations, so limits of
finite-dimensional strategies need not be realized by any finite-dimensional
strategy.  Whether some symmetric game realizes the failure --- a game whose
value is attained by no tensor product strategy, presented together with the
exposing functional as a decision predicate --- is a separate question that
this note does not need to settle.  The local verdict is about the proof: the
cited argument proves the approximate form and does not prove the attainment
form.

\section{Uses of the lemma}

The source invokes the lemma twice.

First, in the soundness proof for the oracularization transformation, the
argument fixes a strategy $\strategy^{\mathrm{orac}}$ for the oracularized game
``with value $1-\eps$'' and applies the lemma to assume without loss of
generality that it is projective.%
\footnote{\path{references/qpbt-paper/10_oracularization.tex},
lines~317--321.}
Here the input is a given strategy, not the value of the game; the
given-strategy form applies verbatim and preserves the value exactly, so the
use is justified by what the proof establishes.

Second, in the soundness proof for answer reduction, the argument assumes the
strict inequality $\val^*(\game) > 1-\eps$ and concludes that there is a
symmetric projective strategy of value greater than $1-\eps$.%
\footnote{\path{references/qpbt-paper/11_answer_reduction.tex},
lines~2322--2332.}
Writing $\val^*(\game) = 1-\eps_0$ with $\eps_0 < \eps$ and applying the
approximate form with $\eps' = (\eps+\eps_0)/2$ produces a symmetric
projective strategy of value at least $1-\eps' > 1-\eps$.  A strict lower
bound on the supremum thus passes through the approximate form; attainment is
not used.

In this project, the soundness statements of the Pauli basis test track
quantify over strategies assumed to succeed with probability at least
$1-\eps$, so any symmetrization step starts from a given strategy and consumes
the given-strategy form.  The blueprint currently records the lemma as part of
the games vocabulary and has no downstream dependency on it; the two source
uses above belong to transformations outside the scope of the present
development.  Consequently, no consumer of the lemma --- in the source or here
--- requires the attainment form.

\section{Comparison with the blueprint}

The blueprint \cite{MIPStarREBlueprint} states the lemma in the source shape.
Its proof records the value-preserving construction and the approximate form
that the construction yields, and explicitly marks the remaining attainment
step as unproved, citing this note; an accompanying remark summarizes the
obstruction.  An earlier revision of the blueprint had instead weakened the
source-labelled conclusion to the approximate form; following the project's
proof-gap protocol, the source-labelled statement is restored and the missing
step is tracked in the proof rather than absorbed into the statement.  No
\Lean{} declaration is currently linked to the lemma.

\section{Conclusion}

The verdict is a proof-level gap in the cited lemma: the symmetrization
argument proves the given-strategy and approximate forms, while the printed
conclusion additionally asserts that the supremum defining $\val^*(\game)$ is
attained, which does not follow.  The gap is harmless for the results built on
the lemma, since both uses in the source, and all planned uses in this
project, proceed from a given strategy of value at least $1-\eps$ or from a
strict lower bound on the value, for which the value-preserving construction
suffices.  For stage~4 (formalization) the recommendation is: formalize the
given-strategy symmetrization lemma as the workhorse statement; keep the
source-labelled lemma in source shape with the attainment step an explicit,
tracked proof obligation; and do not strengthen the source-labelled statement's
hypotheses or weaken its conclusion silently.  If a later investigation
settles the counterexample question for symmetric games, the statement should
then be reclassified from an unproved step to a documented statement
correction.

\bibliographystyle{alpha}
\bibliography{references}

\end{document}
